% TOP_PROBLEM: {"id": 1200005, "title": "Some Questions — Question 5", "category_id": 17, "category": {"id": 17, "name": "group_theory", "display_name": "Group Theory", "description": "Problems about groups, group actions, representations, and related algebraic structures.", "slug": "group-theory", "order_index": 17, "created_at": "2026-07-31T15:26:25.671Z"}, "status": "partially_solved", "problem_number": "AMR-011-0005", "set_id": 13}
% FIRST_POSTED: 2026-10-01
% ENTRY_KIND: statement_only
% UnsolvedMath ID 1200005; source catalogue code AMR-011-0005
% !TeX program = lualatex
% Definition and sources only; this file is not an LLM solution attempt.
\documentclass[11pt,a4paper]{article}
\usepackage[margin=25mm]{geometry}
\usepackage{fontspec}
\setmainfont{lmroman10-regular.otf}[BoldFont=lmroman10-bold.otf,
 ItalicFont=lmroman10-italic.otf,BoldItalicFont=lmroman10-bolditalic.otf]
\usepackage{amsmath,amssymb,mathtools}
\usepackage{unicode-math}
\setmathfont{latinmodern-math.otf}
\AtBeginDocument{\let\setminus\smallsetminus}
\usepackage{xurl,hyperref}
\hypersetup{colorlinks=true,urlcolor=blue}
\setcounter{secnumdepth}{0}
\setlength{\parindent}{0pt}
\setlength{\parskip}{5pt}
\setlength{\emergencystretch}{3em}
\begin{document}
\section*{Some Questions — Question 5}\label{1200005}
{\small UnsolvedMath ID 1200005 · AMR-011-0005 · Group Theory · AMR Open Problem Lists}
\subsection{Definitions and mathematical statement}
Let $C_2$ be the cyclic group of order two. Define finite groups recursively by $W_1=C_2$ and
\[
 W_{n+1}=(W_n\times W_n)\rtimes C_2,
\]
where the nonidentity element of $C_2$ interchanges the two factors. Thus $W_n$ is the full automorphism group of the rooted binary tree of depth $n$, and specifies the permutational iterated wreath product convention unambiguously.

A nontrivial group word is a nonidentity element $w$ of a free group on finitely many variables. Its length $|w|$ is the number of letters in its freely reduced expression, counting inverses as single letters. The word is a law for $W_n$ if every assignment of its variables to elements of $W_n$ evaluates to the identity.

Define
\[
 \ell(n)=\min\{|w|:w\ne1\text{ is a law for }W_n\}.
\]
Determine the asymptotic growth of $\ell(n)$ as $n\to\infty$, with matching upper and lower estimates where possible. The number of variables may vary with the word; the minimization is not restricted to power words or a particular presentation. The power word $x^{2^n}$ is a law and supplies an elementary benchmark, but the problem is whether much shorter identities involving several variables can exist. A word becoming trivial only for selected assignments is not a law. The group parameter is the tree depth, rather than the logarithm of the group order.
\subsection{Short English statement}
How long is the shortest nontrivial identity in the automorphism group of a finite rooted binary tree?
\subsection{Sources}
\begin{itemize}
\item[S1] ULAM.AI and the UnsolvedMath Contributors, supplied problems.json, record 1200005 (AMR-011-0005), AMR Open Problem Lists. Catalogue identity and source transcription; CC BY 4.0.
\url{https://huggingface.co/datasets/ulamai/UnsolvedMath}.
\item[S2] Abert - Some questions (pdf) (2010), Question 5. Original source indexed by AMR-011-0005.
\url{https://www.renyi.hu/~abert/questions.pdf}.
\end{itemize}
\end{document}
